%==============================================================================
\section{Mạng Q sâu (Deep Q-Networks)}
\label{sec:deep-q-networks}
%==============================================================================

\subsection{Deep Q-Network là gì?}

\begin{dinhnghia}[title={Deep Q-Network (DQN)}]
\textbf{Deep Q-Network (DQN)} là thuật toán trong học tăng cường, kết hợp \textbf{mạng nơ-ron sâu} với \textbf{Q-learning}, giúp agent học policy tối ưu trong môi trường phức tạp.
Q-learning truyền thống hiệu quả khi số trạng thái \textbf{nhỏ và hữu hạn}, nhưng gặp khó với không gian trạng thái \textbf{lớn} hoặc \textbf{liên tục} vì kích thước bảng $Q$.
DQN khắc phục bằng cách thay bảng $Q$ bằng mạng nơ-ron \textbf{xấp xỉ} giá trị $Q$ cho mọi cặp trạng thái--hành động.
\end{dinhnghia}

\subsection{Các thành phần chính}

\begin{ynghia}[title={Kiến trúc DQN}]
Dưới đây là các thành phần trong kiến trúc Deep Q-Networks.
\end{ynghia}

\begin{phuongphap}[title={Lớp đầu vào (Input Layer)}]
\textbf{Lớp đầu vào} nhận thông tin trạng thái từ môi trường dưới dạng \textbf{vector giá trị số}.
\end{phuongphap}

\begin{phuongphap}[title={Lớp ẩn (Hidden Layers)}]
\textbf{Lớp ẩn} của DQN gồm nhiều nơ-ron \textbf{fully connected} biến đổi dữ liệu đầu vào thành đặc trưng phức tạp hơn, phù hợp cho \textbf{dự đoán}.
\end{phuongphap}

\begin{phuongphap}[title={Lớp đầu ra (Output Layer)}]
\textbf{Lớp đầu ra}: mỗi hành động khả dĩ ở trạng thái hiện tại tương ứng \textbf{một nơ-ron}.
Giá trị đầu ra của các nơ-ron này là \textbf{ước lượng giá trị} của từng hành động trong trạng thái đó.
\end{phuongphap}

\begin{phuongphap}[title={Bộ nhớ (Memory)}]
\textbf{Bộ nhớ}: DQN dùng \textbf{memory replay} lưu các sự kiện huấn luyện của agent.
Mọi thông tin---trạng thái hiện tại, hành động đã chọn, phần thưởng nhận được, trạng thái kế tiếp---được lưu dưới dạng \textbf{bộ} (tuple) trong bộ nhớ.
\end{phuongphap}

\begin{phuongphap}[title={Hàm mất mát (Loss Function)}]
\textbf{Hàm mất mát}: DQN tính \textbf{chênh lệch} giữa giá trị $Q$ thực tế từ replay memory và giá trị $Q$ dự đoán để xác định loss.
\end{phuongphap}

\begin{phuongphap}[title={Tối ưu hóa (Optimization)}]
\textbf{Tối ưu hóa}: điều chỉnh trọng số mạng để \textbf{tối thiểu hóa} hàm mất mát.
Thường dùng \textbf{stochastic gradient descent (SGD)} cho mục đích này.
\end{phuongphap}

\begin{vidu}[title={Sơ đồ kiến trúc DQN}]
Hình sau minh họa các thành phần trong kiến trúc Deep Q-Network:
\end{vidu}

\tsimg{03_deep_q_networks/img01.jpg}

\subsection{Cách hoạt động của DQN}

\begin{phuongphap}[title={Các bước vận hành}]
Quy trình làm việc của DQN gồm các bước sau.
\end{phuongphap}

\subsubsection{Kiến trúc mạng nơ-ron}

\begin{phuongphap}[title={Neural Network Architecture}]
DQN dùng chuỗi \textbf{frame} (ví dụ ảnh từ game) làm đầu vào và sinh tập giá trị $Q$ cho mọi hành động khả dĩ ở trạng thái đó.
Cấu hình điển hình gồm lớp \textbf{convolutional} nắm quan hệ không gian và lớp \textbf{fully connected} xuất ra giá trị $Q$.
\end{phuongphap}

\subsubsection{Experience Replay}

\begin{phuongphap}[title={Bộ đệm kinh nghiệm}]
Trong huấn luyện, agent lưu các tương tác $(s, a, r, s')$ vào \textbf{replay buffer}.
Lấy mẫu \textbf{ngẫu nhiên} các batch từ buffer để huấn luyện mạng, giảm tương quan giữa kinh nghiệm liên tiếp và tăng \textbf{ổn định} huấn luyện.
\end{phuongphap}

\subsubsection{Target Network}

\begin{phuongphap}[title={Mạng mục tiêu}]
Để ổn định quá trình huấn luyện, DQN dùng \textbf{mạng target riêng} sinh mục tiêu giá trị $Q$.
Mạng target được \textbf{cập nhật định kỳ} trọng số từ mạng chính để giảm rủi ro \textbf{phân kỳ} khi huấn luyện.
\end{phuongphap}

\subsubsection{Chính sách Epsilon-Greedy}

\begin{phuongphap}[title={Epsilon-Greedy Policy}]
Agent dùng chiến lược \textbf{epsilon-greedy}:
\begin{itemize}
  \item Với xác suất $\varepsilon$: chọn hành động \textbf{ngẫu nhiên};
  \item Với xác suất $1-\varepsilon$: chọn hành động có giá trị $Q$ \textbf{cao nhất}.
\end{itemize}
Cân bằng khám phá--khai thác giúp agent học hiệu quả.
\end{phuongphap}

\subsubsection{Quy trình huấn luyện}

\begin{phuongphap}[title={Training Process}]
Mạng nơ-ron được huấn luyện bằng \textbf{gradient descent} để tối thiểu hóa loss giữa $Q$ dự đoán và $Q$ mục tiêu.
Giá trị $Q$ mục tiêu tính theo \textbf{phương trình Bellman}, kết hợp phần thưởng nhận được và giá trị $Q$ \textbf{lớn nhất} của trạng thái kế tiếp.
\end{phuongphap}

\subsection{Hạn chế của Deep Q-Networks}

\begin{luuy}[title={Các hạn chế chính}]
Deep Q-Networks có một số hạn chế ảnh hưởng hiệu quả và hiệu năng:
\begin{enumerate}
  \item \textbf{Bất ổn định} do vấn đề non-stationarity từ việc cập nhật mạng nơ-ron thường xuyên.
  \item Đôi khi \textbf{overestimate} giá trị $Q$, gây tác động tiêu cực lên quá trình học.
  \item Cần \textbf{nhiều mẫu} để học tốt---tốn kém và tốn thời gian tính toán.
  \item Hiệu năng phụ thuộc mạnh vào \textbf{siêu tham số}: learning rate, hệ số chiết khấu, tỷ lệ khám phá.
  \item DQN chủ yếu cho không gian hành động \textbf{rời rạc}; khó với hành động \textbf{liên tục}.
\end{enumerate}
\end{luuy}

\subsection{Double Deep Q-Networks}

\begin{dinhnghia}[title={Double DQN}]
\textbf{Double DQN} là phiên bản mở rộng của DQN, nhằm xử lý vấn đề \textbf{thiên lệch overestimation} trong cập nhật giá trị $Q$.
Thiên lệch này xuất phát từ việc quy tắc cập nhật Q-learning dùng \textbf{cùng một} mạng $Q$ để vừa chọn vừa đánh giá hành động, dẫn tới ước lượng $Q$ \textbf{phồng lên}.
Điều này gây \textbf{bất ổn} huấn luyện và cản trở học.
\end{dinhnghia}

\begin{phuongphap}[title={Hai mạng tách vai trò}]
Double DQN dùng hai mạng:
\begin{itemize}
  \item \textbf{Q-Network}: chọn hành động;
  \item \textbf{Target Network}: đánh giá \textbf{giá trị} của hành động đã chọn.
\end{itemize}
\end{phuongphap}

\begin{phuongphap}[title={Cách tính mục tiêu}]
Thay đổi chính nằm ở cách tính \textbf{target}:
thay vì dùng một mạng $Q$ cho cả chọn và đánh giá hành động kế tiếp, Double DQN dùng mạng $Q$ để \textbf{chọn} hành động ở trạng thái sau và mạng target để \textbf{đánh giá} giá trị $Q$ của hành động đó.
Việc tách này giảm xu hướng overestimate, cho ước lượng giá trị \textbf{chính xác} hơn.
\end{phuongphap}

\begin{tinhchat}[title={Ổn định hơn DQN chuẩn}]
Double DQN mang lại quá trình huấn luyện \textbf{nhất quán và đáng tin cậy} hơn, đặc biệt trong kịch bản như game Atari, nơi DQN thông thường có thể gặp khó với overestimation.
\end{tinhchat}

\subsection{Dueling Deep Q-Networks}

\begin{dinhnghia}[title={Dueling DQN}]
\textbf{Dueling Deep Q-Networks} cải thiện quá trình học của DQN truyền thống bằng cách \textbf{tách} ước lượng giá trị trạng thái khỏi \textbf{advantage} hành động.
Trong DQN cổ điển, mỗi cặp trạng thái--hành động có một giá trị $Q$ riêng biểu diễn phần thưởng tích lũy kỳ vọng.
Cách này có thể \textbf{kém hiệu quả} khi nhiều hành động dẫn tới hậu quả tương tự.
\end{dinhnghia}

\begin{phuongphap}[title={Công thức phân rã $Q$}]
Dueling DQN phân rã giá trị $Q$ thành hai phần:
\begin{itemize}
  \item \textbf{Giá trị trạng thái} $V(s)$: giá trị của việc ở trạng thái $s$;
  \item \textbf{Hàm advantage} $A(s,a)$: mức tốt hơn của hành động $a$ so với các hành động khác trong cùng trạng thái.
\end{itemize}
Giá trị $Q$ được cho bởi:
\[
Q(s,a) = V(s) + A(s,a).
\]
\end{phuongphap}

\begin{tinhchat}[title={Lợi ích học tách rời}]
Dueling DQN giúp agent hiểu môi trường sâu hơn và tránh học các ước lượng giá trị hành động \textbf{không cần thiết}.
Hiệu năng cải thiện, đặc biệt khi phần thưởng \textbf{trễ}, vì agent nắm rõ hơn tầm quan trọng của các trạng thái khi chọn hành động tối ưu.
\end{tinhchat}
